\scanpage{297}

% letterhead
\begin{center}
U N I V E R S I T E\\
DES SCIENCES ET TECHNIQUES\\
DU LANGUEDOC\\
\texttt{------------}
\end{center}

MONTPELLIER, le 26 Juin 1980

Service de la Scolarité\\
MR/GJ\\
Diffusion : Liste H\\
N°: 3713

% printed form
\uline{AVIS DE SOUTENANCE D'UNE THESE DE DOCTORAT D'ETAT}

\uline{Nom et prénom du candidat} :TOUBOUL André

\uline{Titre de la thèse} : ``Caractérisation d'un dispositif à transfert de charges (S.C.C.D.) Transfert et bruit de fond par traitement numérique''.

\uline{Date, heure et lieu de soutenance prévus} : 11 JUILLET 1980 à 10 H. 30 mn\\
Salle de cours n° 9\\
(Bâtiment Administratif)

\uline{Nom des Directeurs de thèse} : Professeur G. LECOY et A.A. WALMA

\uline{Désignation du laboratoire où la thèse est préparée} : G.E.M.C.E.S.-C.E.E.S. (U.S.T.L.)

\uline{Résumé de la thèse} : Après avoir rappelé le fonctionnement d'un dispositif à transfert de charges (C.C.D.) à canal en surface, on présente une méthode numérique originale de mesure du bruit de fond menant directement à la variance des charges.Des mesures de bruit sont données : elles montrent la prépondérance des mécanismes de transfert sur le bruit. Pour certains modes de fonctionnement, le produit de l'inefficacité de transfert par la taille du paquet peut rendre compte de la variance des charges de bruit. L'influence des modes d'injection des charges ainsi que du taux de remplissage des puits a été mise en évidence. Pour expliquer le comportement du bruit et de l'inefficacité de transfert, on a suggéré que le rôle de la structure amorphe de l'interface soit considéré.

\uline{Durée approximative des recherches ayant abouti à la thèse} : 6 ans

\uline{Nom des rapporteurs} :

M. SAVELLI M., Professeur, E.E.A., U.S.T.L.

M. LECOY G., Professeur, E.E.A., U.S.T.L.

\unsure{M.|Mme} WALMA Alle Anne, Professeur Associé, E.E.A., U.S.T.L.

M. HOBSON G.S., Professeur, Electrical Eng., Université de Sheffield (G.B.)

\uline{Composition du jury proposé} :

M. SAVELLI, E.E.A., Professeur, U.S.T.L.  (Président)

M. BOREL, Directeur Recherche et Développement EFCIS, GRENOBLE

M. CAMUS , Directeur, CNET (Meylan-Grenoble)

M. DUMONTET, Phys. Sol. \& Phys. At, Professeur, U.S.T.L.\\
(Membre étranger à la spécialité du candidat)

M. HOBSON, Electrical Eng., Professeur, Univ. Sheffield (G.B.)\\
(Membre extérieur choisi en raison de sa compétence scientifique)

M. NOUGIER, E.E.A, Professeur, U.S.T.L.

M. LECOY, E.E.A., Professeur, U.S.T.L.

M. WALMA ,E.E.A., Professeur associé, U.S.T.L.

P. le Secrétaire Général,\\
L'Attachée d'Administration Universitaire,

\handwritten{\unsure{Gentil}}

\uline{A. GENTIL}

% archive sheet number: 296
